\documentclass[a4paper,12pt]{article}
\usepackage[T2A]{fontenc}
\usepackage[utf8]{inputenc}
\usepackage[russian]{babel}
\usepackage{amsmath,amssymb,mathtools}
\usepackage{PTSans}
\renewcommand{\familydefault}{\sfdefault}
\usepackage{icomma}
\usepackage[margin=20mm,headheight=15pt,headsep=7mm,footskip=12mm]{geometry}
\usepackage{microtype}
\usepackage{parskip}
\usepackage{enumitem}
\usepackage{booktabs,tabularx,array}
\usepackage{tikz}
\usetikzlibrary{arrows.meta,positioning,shapes.geometric,calc}
\usepackage{pgfplots}
\pgfplotsset{compat=1.18}
\usepackage{xcolor}
\usepackage{fancyhdr}
\usepackage{hyperref}
\definecolor{accent}{HTML}{008080}
\definecolor{ink}{HTML}{282F33}
\definecolor{muted}{HTML}{616C70}
\definecolor{pale}{HTML}{E2EEEE}
\hypersetup{colorlinks=true,linkcolor=accent,urlcolor=accent}
\color{ink}
\setlength{\parindent}{0pt}
\setlength{\parskip}{5pt}
\setlist[itemize]{leftmargin=*,topsep=3pt,itemsep=2pt,parsep=0pt}
\setlist[enumerate]{leftmargin=*,topsep=3pt,itemsep=4pt,parsep=0pt}
\renewcommand{\arraystretch}{1.25}
\newcommand{\key}[1]{\textcolor{accent}{\textbf{#1}}}
\newcommand{\smallnote}[1]{{\small\textcolor{muted}{#1}}}
\newcommand{\topic}[1]{\vspace{5pt}{\color{accent}\large\bfseries #1}\par\vspace{2pt}}
\pagestyle{fancy}
\fancyhf{}
\fancyhead[L]{\small\textcolor{muted}{Прикладные задачи ЕГЭ}}
\fancyhead[R]{\small\textcolor{muted}{Анна Трапезникова}}
\fancyfoot[C]{\small\textcolor{muted}{\thepage}}
\renewcommand{\headrulewidth}{0pt}
\renewcommand{\footrulewidth}{0pt}
\begin{document}
\thispagestyle{empty}
\begin{center}
\vspace*{29mm}
{\large\textcolor{accent}{МАТЕМАТИКА}}\par
\vspace{20mm}
{\Huge\mdseries Чек-лист}\par
\vspace{14mm}
{\LARGE\textcolor{accent}{Задачи прикладного характера}}\par
\vspace{4mm}
{\Large\textcolor{accent}{Формулы, уравнения и смысл ответа}}\par
\vspace{18mm}
{\normalsize Подготовка к ЕГЭ по профильной математике}\par
\vfill
{\large Анна Трапезникова}\par
\vspace{8mm}
{27 октября 2026 г.}\par
\vspace{15mm}
{\small Лёвин Артём Александрович}\par
{\small эксклюзивно для Анны Трапезниковой}\par
\vspace{14mm}
\begin{tikzpicture}[x=1cm,y=1cm]
  \draw[accent,line width=.75pt] (0,0) .. controls (1.4,.62) and (2.1,-.24) .. (3.5,.05)
    .. controls (4.7,.49) and (5.45,-.24) .. (6.75,.02)
    .. controls (7.9,.39) and (8.65,-.43) .. (10,0);
  \fill[accent] (4.65,.08) circle (1.2pt);
\end{tikzpicture}
\vspace*{8mm}
\end{center}
\clearpage

\section*{1. Что нужно увидеть в условии}
\key{Главная идея.} За длинным прикладным условием обычно скрывается знакомое линейное, квадратное или простейшее степенное уравнение. Основная работа --- правильно перевести текст задачи на математический язык и осмысленно выбрать ответ.

\begin{enumerate}[label=\textcolor{accent}{\arabic*.}]
\item \textbf{Определите искомое.} Время, высота, цена, скорость, радиус или длительность --- разные величины, даже если в конце они связаны одним уравнением.
\item \textbf{Выделите исходные данные.} Выпишите буквенные обозначения, числа и требуемые ограничения.
\item \textbf{При необходимости нарисуйте схему.} Особенно полезно для воды в баке, высоты над стеной, полёта тела, нагревания.
\item \textbf{Согласуйте единицы измерения.} Метры и миллиметры нельзя подставлять в одну формулу как одинаковые единицы.
\item \textbf{Подставьте данные и решите уравнение.} Если формул две, одну часто удобно подставить в другую.
\item \textbf{Проверьте смысл корней.} Отрицательная длина, момент после отключения прибора и время после опустошения бака не подходят.
\end{enumerate}

\topic{Ключевые слова задачи}
\begin{tabularx}{\linewidth}{@{}p{.31\linewidth}X@{}}
\toprule
\textbf{Формулировка} & \textbf{Что означает математически}\\
\midrule
«Увеличилось на $d$» & Новое значение минус первоначальное равно $d$.\\
«Не менее $A$» & Искомая величина $\ge A$.\\
«Не более $A$» & Искомая величина $\le A$.\\
«Наименьшее значение» & Из допустимых значений выбрать минимальное.\\
«Максимальное значение» & Из допустимых значений выбрать максимальное.\\
«В течение какого времени» & Найти длину временного промежутка: $\Delta t=t_2-t_1$.\\
\bottomrule
\end{tabularx}

\smallnote{Важно: заменять неравенство равенством разрешено для нахождения его граничных значений; затем необходимо проверить направление неравенства и смысл задачи.}

\section*{2. Подготовка вычислений}

\topic{Перевод единиц и степени десяти}
\[
1\ \text{мм}=10^{-3}\ \text{м},\qquad
4{,}5\ \text{мм}=4{,}5\cdot10^{-3}\ \text{м}=45\cdot10^{-4}\ \text{м}.
\]
\[
1{,}2\cdot10^{-5}=12\cdot10^{-6},\qquad
\frac{10^{-4}}{10^{-5}}=10^{1}=10.
\]
\key{Проверка:} при делении степеней с основанием $10$ показатели \emph{вычитаются}: $10^a:10^b=10^{a-b}$.

\clearpage
\section*{3. Линейные модели: изменение, прибыль, уровень воды}
\topic{Удлинение рельса}
При температуре $0^\circ\!\mathrm{C}$ длина рельса $l_0=10$ м. Коэффициент теплового расширения $\alpha=1{,}2\cdot10^{-5}$, а удлинение равно $4{,}5$ мм. Формула:
\[
l(t)=l_0(1+\alpha t).
\]
Раскрываем скобки и выделяем именно \emph{изменение длины}:
\[
l-l_0=l_0\alpha t,\qquad
45\cdot10^{-4}=10\cdot12\cdot10^{-6}\cdot t.
\]
Отсюда
\[
t=\frac{45\cdot10^{-4}}{120\cdot10^{-6}}
 =\frac{45}{12}\cdot10=37{,}5.
\]
\key{Ответ:} $37{,}5^\circ\!\mathrm{C}$.

\topic{Минимальный выпуск продукции}
Прибыль предприятия моделируется формулой
\[
\Pi=q(p-v)-F,
\]
где $q$ --- количество изделий, $p$ --- цена, $v$ --- затраты на одно изделие, $F$ --- постоянные расходы.

Если $p=600$, $v=300$, $F=700\,000$ и прибыль должна быть \emph{не менее} $500\,000$, то
\[
300q-700\,000\ge500\,000
\ \Longrightarrow\ 300q\ge1\,200\,000
\ \Longrightarrow\ q\ge4000.
\]
\key{Ответ:} минимальный выпуск --- $4000$ изделий. Здесь переход к границе допустим, поскольку прибыль возрастает при увеличении $q$: $p-v=300>0$.

\topic{Колодец: изменение высоты}
Расстояние от края колодца до поверхности воды находится по формуле $h(t)=5t^2$ (метры, время падения камешка в секундах). До дождя камешек падал $0{,}6$ с, после дождя --- $0{,}4$ с. Уровень воды поднялся на разность глубин:
\[
\Delta h=h(0{,}6)-h(0{,}4)
=5\bigl(0{,}6^2-0{,}4^2\bigr)=5\cdot0{,}2=1\ \text{м}.
\]
\key{Приём:} $a^2-b^2=(a-b)(a+b)$. Разность уровней обязательно вычисляют по \emph{двум} моментам.

\clearpage
\section*{4. Две формулы и физический смысл}
\topic{Спрос и выручка: подстановка}
Пусть спрос $q=60-5p$, выручка $r=qp$, причём $r=160$ (в тысячах рублей). Тогда
\[
(60-5p)p=160\ \Longrightarrow\
 p^2-12p+32=0.
\]
\[
(p-4)(p-8)=0\ \Longrightarrow\ p=4\ \text{или}\ p=8.
\]
Если спрашивают \emph{максимальную цену}, выбирают $p=8$; обе цены допустимы, поскольку $q\ge0$.

\topic{Ведро на вертикальной окружности}
В верхней точке траектории силу давления на дно описывает
\[
P=m\left(\frac{v^2}{L}-g\right).
\]
Если вода не должна выливаться, на границе допустимой скорости $P=0$. Масса $m>0$, поэтому
\[
\frac{v^2}{L}-g=0\ \Longrightarrow\ v^2=gL.
\]
При $L=0{,}625$ м и $g=10$ м/с$^2$:
\[
v=\sqrt{10\cdot0{,}625}=\sqrt{6{,}25}=2{,}5\ \text{м/с}.
\]
\key{Ответ:} минимальная скорость в верхней точке --- $2{,}5$ м/с. Скорость берётся неотрицательной.

\topic{Архимедова сила: куб неизвестной}
Если $F_A=\alpha\rho g r^3$, то радиус находится из
\[
r^3=\frac{F_A}{\alpha\rho g},\qquad
r=\sqrt[3]{\frac{F_A}{\alpha\rho g}}.
\]
При $F_A=336\,000$ Н, $\alpha=4{,}2$, $\rho=1000$ кг/м$^3$ и $g=10$ м/с$^2$ получаем $r^3=8$, то есть $r=2$ м.

\topic{Текстовая ловушка: пролететь выше стены}
Если высота стены равна $9$ м, а камень должен пройти минимум на $1$ м \emph{выше} неё, требуемая высота камня --- $9+1=10$ м. В модели $y=ax^2+bx$ на расстоянии до стены $x=d$ нужно проверить $y(d)\ge10$.

\smallnote{В уравнениях с $v^2$, $r^3$ и другими степенями учитывайте, какой знак допускает реальная величина.}

\clearpage
\section*{5. Квадратная модель: какое время нужно?}
\topic{Мяч выше заданной высоты}
Высота мяча задаётся формулой
\[
h(t)=1{,}4+9t-5t^2.
\]
Границы пребывания \emph{не ниже $3$ м} находятся из уравнения
\[
1{,}4+9t-5t^2=3\ \Longleftrightarrow\ 5t^2-9t+1{,}6=0.
\]
\[
D=81-32=49,\quad t_1=\frac{9-7}{10}=0{,}2,\quad
 t_2=\frac{9+7}{10}=1{,}6.
\]
\key{Ищут длительность}, поэтому вычитаем:
\[
\boxed{\Delta t=t_2-t_1=1{,}6-0{,}2=1{,}4\ \text{с}.}
\]
\begin{center}
\begin{tikzpicture}
\begin{axis}[
width=.83\linewidth,height=6.2cm,
xmin=0,xmax=1.95,ymin=0,ymax=6,
axis lines=middle, xlabel={$t$, с},ylabel={$h$, м},
xtick={0,.2,.9,1.6},xticklabels={$0$,$0{,}2$,$0{,}9$,$1{,}6$},
ytick={0,3,5.45},yticklabels={$0$,$3$,$5{,}45$},
clip=false,grid=both,minor tick num=0,
major grid style={draw=pale,thin},tick label style={font=\scriptsize},
label style={font=\small},every axis plot/.append style={smooth}]
\addplot[accent,very thick,domain=0:1.94,samples=201] {1.4+9*x-5*x^2};
\addplot[muted,densely dashed,domain=0:1.92] {3};
\addplot[only marks,mark=*,mark size=2pt,accent] coordinates {(.2,3) (1.6,3)};
\draw[<->,>=Latex,accent,thick] (axis cs:.2,3.34)--(axis cs:1.6,3.34)
 node[midway,above,font=\small] {$\Delta t=1{,}4$ с};
\end{axis}
\end{tikzpicture}
\end{center}
\smallnote{Промежуток выше уровня $3$ м расположен между двумя моментами пересечения графика с прямой $h=3$.}

\topic{Прибор отключается при достижении температуры}
Пусть $T(t)=1400+200t-10t^2$. При достижении $1760$ К нагреватель \emph{автоматически отключается}. Решаем уравнение:
\[
1760=1400+200t-10t^2\ \Longleftrightarrow\
 t^2-20t+36=0.
\]
Получаем $t=2$ и $t=18$. \key{Ответ: $2$ минуты}, потому что первое достижение порога приводит к отключению. Момент $18$ минут относится к продолжению математической модели уже после выключения.

\clearpage
\section*{6. Квадратная модель: проверить границы процесса}
\topic{Вода вытекает из цилиндрического бака}
Для высоты воды используется модель
\[
H(t)=H_0-\sqrt{2gH_0}\,kt+\frac{g}{2}k^2t^2.
\]
Если бак цилиндрический (с постоянной площадью сечения), четверть первоначального \emph{объёма} соответствует высоте $H_0/4$. При $H_0=20$, $g=10$, $k=1/500$ получаем
\[
5=20-\frac{t}{25}+\frac{t^2}{50\,000}.
\]
Умножаем на $50\,000$:
\[
t^2-2000t+750\,000=0,
\qquad t=500\ \text{или}\ t=1500.
\]
\key{Ответ: $500$ секунд.} Вода полностью заканчивается при $t=1000$ с (из $H(t)=0$). Поэтому значение $1500$ с уже не относится к физическому процессу.

\topic{Момент инерции: смысл корней}
Для катушки используется формула
\[
I=\frac{(m+2M)R^2}{2}+M(2Rh+h^2).
\]
Пусть $m=8$, $M=1$, $R=10$, $I=625$. После подстановки:
\[
625=500+20h+h^2\ \Longleftrightarrow\
 h^2+20h-125=0.
\]
\[
h=5\ \text{или}\ h=-25.
\]
\key{Ответ: $5$ см.} Отрицательная высота в данном условии лишена физического смысла.

\topic{Когда температура «не менее» заданной}
Если $T(t)=-t^2+14t+25$ и требуется время, в течение которого температура \emph{не ниже} $70^\circ$C, составляем неравенство:
\[
-t^2+14t+25\ge70\ \Longleftrightarrow\
 (t-5)(t-9)\le0.
\]
Следовательно, $5\le t\le9$, а \emph{длительность} равна
\[
\boxed{9-5=4\ \text{мин}.}
\]
\key{Не перепутайте:} первый момент достижения порога --- $5$ минут; длительность пребывания над порогом --- $4$ минуты.

\topic{Что обязательно проверить}
\begin{itemize}
\item $t\ge0$; после отключения или опустошения модель может потерять физический смысл.
\item Длина, расстояние, радиус и масса не могут оказаться отрицательными.
\item «Не менее» и «не более» требуют проверки промежутка и его границ.
\item Если найдены два момента, уточните, спрашивают ли \emph{один момент} или \emph{продолжительность}.
\end{itemize}

\clearpage
\section*{Алгоритм решения прикладной задачи}
\vspace{11mm}
\begin{center}
\begin{tikzpicture}[>=Latex,node distance=8mm,
act/.style={draw=accent!68,rounded corners=3pt,align=center,text width=11cm,minimum height=1.2cm,font=\normalsize,inner sep=6pt},
terminal/.style={draw=accent,ellipse,align=center,text width=5.4cm,minimum height=.85cm,font=\normalsize,inner sep=3pt},
choice/.style={draw=accent!75,diamond,aspect=3.3,align=center,text width=5cm,inner sep=3pt,font=\small},
arr/.style={->,very thick,draw=muted}]
\node[terminal] (start) {Прочитайте задачу};
\node[act,below=of start] (values) {Обозначьте искомое и выпишите данные\\Уточните, что значит «на», «не менее», «максимальное»};
\node[act,below=of values] (pic) {Если нужен физический рисунок, сделайте его\\Приведите величины к согласованным единицам};
\node[act,below=of pic] (calc) {Подставьте данные в формулу\\Решите уравнение или неравенство};
\node[choice,below=12mm of calc] (interval) {Нужна длительность\ времени?};
\node[act,text width=5.1cm,below left=17mm and 5mm of interval] (length) {Найдите допустимый интервал\\Вычислите $t_2-t_1$};
\node[act,text width=5.1cm,below right=17mm and 5mm of interval] (moment) {Выберите допустимый момент или значение по условию};
\node[act,below=33mm of interval] (check) {Проверьте единицы, ограничения и смысл ответа};
\node[terminal,below=of check] (end) {Запишите ответ};
\draw[arr] (start)--(values);
\draw[arr] (values)--(pic);
\draw[arr] (pic)--(calc);
\draw[arr] (calc)--(interval);
\draw[arr] (interval.south west)-- node[above left,font=\small,muted]{да} (length.north);
\draw[arr] (interval.south east)-- node[above right,font=\small,muted]{нет} (moment.north);
\draw[arr] (length.south) |- (check.west);
\draw[arr] (moment.south) |- (check.east);
\draw[arr] (check)--(end);
\end{tikzpicture}
\end{center}

\clearpage
\section*{Тренировочный блок --- Путь А}
\smallnote{10 задач по нарастающей сложности. Решения выполните самостоятельно. Ответы приведены после заданий.}
\vspace{4mm}
\begin{enumerate}[label=\textcolor{accent}{\textbf{\arabic*.}},leftmargin=*,itemsep=13pt,parsep=4pt]
\item При температуре $0^\circ\mathrm{C}$ рельс имеет длину $l_0=12$ м. При возрастании температуры происходит тепловое расширение рельса, и его длина, выраженная в метрах, меняется по закону
\[
 l(t^\circ)=l_0(1+\alpha\cdot t^\circ),
\]
где $\alpha=1{,}5\cdot10^{-5}\ ({}^\circ\mathrm{C})^{-1}$ --- коэффициент теплового расширения, $t^\circ$ --- температура (в градусах Цельсия). При какой температуре рельс удлинится на $5{,}4$ мм? Ответ выразите в градусах Цельсия.

\item После дождя уровень воды в колодце может повыситься. Мальчик измеряет время падения $t$ небольших камешков в колодец и рассчитывает расстояние до воды по формуле
\[
 h=5t^2,
\]
где $h$ --- расстояние в метрах, $t$ --- время падения в секундах, прошедшее с момента броска. До дождя время падения камешков составляло $1{,}2$ с. На сколько должен подняться уровень воды после дождя, чтобы измеряемое время изменилось на $0{,}4$ с? Ответ выразите в метрах.

\item Некоторая компания продаёт свою продукцию по цене $p=750$ руб. за единицу, переменные затраты на производство одной единицы продукции составляют $v=450$ руб., постоянные расходы предприятия $f=480\,000$ руб. в месяц. Месячная операционная прибыль предприятия, выраженная в рублях, вычисляется по формуле:
\[
 \pi(q)=q(p-v)-f.
\]
Определите наименьший месячный объём производства $q$ (единиц продукции), при котором месячная операционная прибыль предприятия будет не менее $420\,000$ руб.

\item Зависимость объёма спроса $q$ на продукцию предприятия-монополиста от цены $p$ (тыс. руб.) задаётся формулой:
\[
 q=50-2p.
\]
Выручка предприятия за месяц $r$ (в тыс. руб.) определяется как $r(p)=q\cdot p$. Определите максимальный уровень цены $p$, при котором месячная выручка $r(p)$ составит не менее $288$ тыс. руб. Ответ приведите в тыс. руб.

\newpage\item Если достаточно быстро вращать ведёрко с водой на верёвке в вертикальной плоскости, то вода не будет выливаться. При вращении ведёрка сила давления воды на дно не остаётся постоянной: она максимальна в нижней точке и минимальна в верхней. Вода не будет выливаться, если сила её давления на дно будет положительной во всех точках траектории, кроме верхней, где она может быть равной нулю. В верхней точке сила давления, выраженная в ньютонах, равна
\[
 P=m\left(\frac{v^2}{L}-g\right),
\]
где $m$ --- масса воды в килограммах, $v$ --- скорость движения ведёрка в м/с, $L$ --- длина верёвки в метрах, $g$ --- ускорение свободного падения (считайте $g=10$ м/с$^2$). С какой минимальной скоростью надо вращать ведёрко, чтобы вода не выливалась, если длина верёвки равна $0{,}9$ м? Ответ выразите в м/с.

\item На верфи инженеры проектируют новый аппарат для погружения на небольшие глубины. Конструкция имеет форму сферы, а значит, действующая на аппарат выталкивающая (архимедова) сила, выражаемая в ньютонах, будет определяться по формуле:
\[
 F_A=\alpha\rho gr^3,
\]
где $\alpha=3$ --- постоянная, $r$ --- радиус аппарата в метрах, $\rho=1000$ кг/м$^3$ --- плотность воды, а $g$ --- ускорение свободного падения (считайте $g=10$ Н/кг). Каков может быть максимальный радиус аппарата, чтобы выталкивающая сила при погружении была не больше, чем $240\,000$ Н? Ответ выразите в метрах.

\item Деталью некоторого прибора является вращающаяся катушка. Она состоит из трёх однородных соосных цилиндров: центрального массой $m=6$ кг и радиуса $R=6$ см, и двух боковых с массами $M=2$ кг и радиусами $R+h$. При этом момент инерции катушки относительно оси вращения, выражаемый в кг$\cdot$см$^2$, даётся формулой
\[
 I=\frac{(m+2M)R^2}{2}+M(2Rh+h^2).
\]
При каком максимальном значении $h$ момент инерции катушки не превысит предельного значения $270$ кг$\cdot$см$^2$? Ответ выразите в сантиметрах.

\newpage\item Высота над землёй подброшенного вверх мяча меняется по закону
\[
 h(t)=1+10t-2t^2,
\]
где $h$ --- высота в метрах, $t$ --- время в секундах, прошедшее с момента броска. Сколько секунд мяч будет находиться на высоте не менее девяти метров?

\item В боковой стенке высокого цилиндрического бака и у самого дна закреплён кран. После его открытия вода начинает вытекать из бака, при этом высота столба воды в нём, выраженная в метрах, меняется по закону
\[
 H(t)=H_0-\sqrt{2gH_0}\cdot kt+\frac{g}{2}k^2t^2,
\]
где $t$ --- время (в секундах), прошедшее с момента открытия крана, $H_0=20$ м --- начальная высота столба воды, $k=\dfrac1{250}$ --- отношение площадей поперечных сечений крана и бака, а $g$ --- ускорение свободного падения (считайте $g=10$ м/с$^2$). Через сколько секунд после открытия крана в баке останется четверть первоначального объёма?

\item Электрический прибор после включения в сеть нагревается и через некоторое время автоматически выключается и начинает остывать. Температура прибора в градусах Цельсия вычисляется по формуле $T=-t^2+16t+T_0$, где $t$ --- время в минутах, прошедшее после включения прибора в сеть, $T_0=20^\circ\mathrm{C}$ --- начальная температура прибора. На протяжении какого времени в минутах температура прибора будет не менее $80^\circ\mathrm{C}$?
\end{enumerate}

\clearpage
\section*{Ответы к тренировочному блоку}
\vspace{8mm}
\begin{center}
\begin{tabularx}{.88\linewidth}{@{}cX@{}}
\toprule
\textbf{№} & \textbf{Ответ}\\
\midrule
1 & $30^\circ\mathrm{C}$\\
2 & $4$ м\\
3 & $3000$ единиц продукции\\
4 & $16$ тыс. руб.\\
5 & $3$ м/с\\
6 & $2$ м\\
7 & $3$ см\\
8 & $3$ с\\
9 & $250$ с\\
10 & $4$ мин\\
\bottomrule
\end{tabularx}
\end{center}
\vspace{13mm}
\topic{Короткая проверка перед сдачей}
\begin{itemize}
\item Величина в ответе совпадает с той, о которой спрашивают?
\item Если спрашивали «сколько времени», найдена ли разность моментов?
\item Если корней несколько, проверены ли ограничения и физический смысл?
\item Если нужно «не менее» или «не более», выбран ли правильный интервал?
\item Выполнены ли перевод единиц и округление, если они требуются?
\end{itemize}
\vfill
\begin{center}
{\small\textcolor{muted}{Лёвин Артём Александрович эксклюзивно для Анны Трапезниковой}}
\end{center}
\end{document}